% =====================================================================
%  Samle-mappe: elektromekaniske interaksjoner mellom vindturbinen
%  (IEA 15 MW) og LEOGO-plattformnettet.
%
%  Denne fila (00_overview.tex) er en INDEKS/leserekkefølge. Den samler de
%  nye seksjonene som ble skrevet i etterkant av overleveringsnotatet
%  (docs/STATUS_OG_VEIVALG.md), og knytter dem til de eksisterende
%  seksjonene i ../ (docs/).
%
%  Bruk: enten \input hver enkelt fil i hoveddokumentet i rekkefølgen under,
%  eller \input denne fila direkte (den \input-er resten).
%
%  Preamble (kreves av alle filene her):
%    \usepackage{amsmath, amssymb}
%    \usepackage{graphicx}
%    \usepackage{booktabs}
%    \usepackage{siunitx}
%    \DeclareSIUnit{\pu}{pu}   % per-unit, brukt av \SI{}{\pu} i disse filene
%    \usepackage{hyperref}     % for kilde-lenkene (\href) i Sources-blokkene
%
%  Figurstier: figurene ligger i ../results/tower_test/ (PNG). Kopier dem til
%  ditt \graphicspath, eller juster \includegraphics-stiene i hver fil.
% =====================================================================

% ---------------------------------------------------------------------
%  Leserekkefølge (elektromekanisk interaksjon)
% ---------------------------------------------------------------------
%
%  Eksisterende seksjoner (i ../, uendret):
%    1. frequency_analysis.tex          \label{sec:frequency-analysis}
%    2. forced_response.tex             \label{sec:forcedresponse}
%    3. drivetrain_torsion_resonance.tex\label{sec:fmu-torsion-resonance}
%    4. openfast_excitation.tex         \label{sec:openfast-excitation}
%    5. network_excitation.tex          \label{sec:network-excitation}
%    6. tower_ss_resonance.tex          \label{sec:tower-ss-resonance}
%
%  Nye seksjoner (denne mappen), i anbefalt rekkefølge:
%    0. narrative.tex                   \label{sec:narrative}   (rød tråd, les først)
%    7. simplified_tower_model.tex      \label{sec:simplified-tower-model}
%    8. rosco_blocking.tex              \label{sec:rosco-blocking}
%    9. shaft_excitation.tex            \label{sec:shaft-excitation}
%   10. outer_droop.tex                 \label{sec:outer-droop}
%   11. outer_droop_necessity.tex       \label{sec:outer-droop-necessity}
%   12. new_tests.tex                   \label{sec:new-tests}
%
%  De nye seksjonene bygger på den REDUSERTE (FMU-frie) modellen
%  WindTurbineTower, som samler drivverk + tårn-SS + tårn-fore-aft i én
%  billig TOPS-modell. Der de eksisterende FMU-seksjonene viser at koblingen
%  finnes i høy-fidelitet, viser de nye seksjonene at den kan drives, styres
%  og resonans-testes systematisk i sanntid mot hele LEOGO-nettet.
%
%  narrative.tex er en kort kronologisk "rød tråd" (4 faser) til poster/pitch:
%  Fase 1 = droop-kontroll + laststeg i LEOGO-nettet (starten); Fase 4 lukker
%  sirkelen tilbake til droopen (frekvensstøtte-/tårnlast-avveiningen). Les først.
% ---------------------------------------------------------------------

% =====================================================================
%  Seksjon: Rød tråd / narrativ for hele sommerprosjektet.
%  Formål:  Kort kronologisk fortelling (4 faser) til poster/pitch.
%           Fase 1 = droop-kontroll + laststeg i LEOGO-nettet (STARTEN),
%           Fase 4 lukker sirkelen tilbake til droopen.
%  Kilder:  casestudies/dyn_sim/test_WT_LEOGO_droop_comparison_diagnostics.py
%           casestudies/dyn_sim/test_WT_LEOGO_torsional_resonance_sim.py
%           casestudies/dyn_sim/sweep_resonance.py            (OL GenTq)
%           casestudies/dyn_sim/test_WT_LEOGO_tower_sim.py
%           src/tops_openfast/dyn_models/windturbine_tower.py
%  Preamble: amsmath, amssymb, graphicx, booktabs, siunitx, hyperref
%            \DeclareSIUnit{\pu}{pu}
% =====================================================================

\section{Red Thread: From Grid Frequency Support to Tower Loading}
\label{sec:narrative}

This project studies electro-mechanical interactions in an offshore wind
turbine --- the IEA \SI{15}{\mega\watt} reference turbine --- coupled through a
grid-forming converter to the LEOGO islanded oil-and-gas platform network, and
co-simulated in TOPS. The study is best read as a single storyline in four
phases; each phase is developed in full in the section it points to. The thread
runs from a grid \emph{service} the turbine can provide to a structural
\emph{cost} that same service imposes, and back again.

\paragraph{Phase 1 --- Frequency-support droop and grid load steps (the starting
point).}
The turbine is first given an explicit frequency-support droop and de-loaded to
hold a headroom reserve. Matched cases \emph{with} and \emph{without} the droop
are then run on the LEOGO grid under load steps applied both abruptly and
gradually, and the droop is shown to reduce the grid-frequency excursion (the
nadir) after the event. This establishes the operational motivation for
everything that follows: on an islanded platform the wind turbine \emph{can} help
arrest frequency deviations, the same primary-reserve service the gas-turbine
governors provide. See \S\ref{sec:outer-droop} and \S\ref{sec:outer-droop-necessity};
the matched droop/no-droop experiment is
\texttt{casestudies/dyn\_sim/test\_WT\_LEOGO\_droop\_comparison\_diagnostics.py}.

\paragraph{Phase 2 --- Modal analysis and the shaft (drivetrain-torsion) mode.}
A modal analysis of the coupled turbine--LEOGO system identifies a lightly damped
drivetrain torsion (shaft) mode, and forced-response tests confirm that an
\emph{electrical} disturbance can drive it into resonance. The mode is excited in
both models --- the reduced TOPS drivetrain (a sharp resonance under a shunt-load
sweep) and the OpenFAST FMU --- and, crucially, it is reachable from the LEOGO
network bus itself, though attenuated by the network impedance. This is the first
concrete electrical-to-mechanical resonance of the study. See
\S\ref{sec:frequency-analysis}, \S\ref{sec:forcedresponse},
\S\ref{sec:fmu-torsion-resonance} and \S\ref{sec:network-excitation}.

\paragraph{Phase 3 --- Tower modes in OpenFAST: ROSCO blocking and open-loop
injection.}
Turning to the tower, the side-to-side (SS) and fore-aft (FA) modes are examined
in the high-fidelity OpenFAST turbine. Here the normal electrical path is found
to be \emph{blocked} by the ROSCO controller, which computes the generator torque
internally and rejects the external commands. Injecting torque instead through
ROSCO's open-loop input file excites the SS mode --- a sharp resonance peaking in
the tower band --- but \emph{not} the FA mode, because the FA mode is thrust
driven rather than torque driven. This establishes that the electro-mechanical
coupling is frequency- and mode-selective: the generator-torque path reaches SS,
not FA. See \S\ref{sec:openfast-excitation}, \S\ref{sec:rosco-blocking} and
\S\ref{sec:tower-ss-resonance}.

\paragraph{Phase 4 --- Tower modes in the reduced model, operational scenarios,
and closing the circle.}
Because ROSCO blocks the genuinely closed electrical loop in the FMU, the SS and
FA tower modes are added as calibrated modal oscillators to the reduced
\texttt{WindTurbineTower} model, which has an \emph{open} electrical-to-mechanical
path and runs in real time against the full LEOGO grid. Operational scenarios
(gas-turbine trip, a resonant cyclic process load, and a frequency sweep) then
close the loop back to Phase~1: the very frequency-support droop that improved the
grid frequency is shown to ring the SS tower mode \emph{harder}. That is the
central result --- the frequency-support / tower-loading trade-off. See
\S\ref{sec:simplified-tower-model}, \S\ref{sec:shaft-excitation} and
\S\ref{sec:new-tests}, in particular the trade-off in \S\ref{subsec:nt-tradeoff}.

\paragraph{The thread in one sentence.}
A wind turbine that supports an islanded platform grid trades grid-frequency
quality against structural tower loading, because the same generator-torque path
that delivers fast active power also drives the lightly damped side-to-side tower
mode.

% =====================================================================
%  Seksjon: Implementering av tårnmodene (side-to-side + fore-aft) i den
%           forenklede vindturbinmodellen (WindTurbineTower).
%  Kilde:   src/tops_openfast/dyn_models/windturbine_tower.py
%  Preamble: amsmath, amssymb, graphicx, booktabs, siunitx, hyperref
%           + \DeclareSIUnit{\pu}{pu}   (per-unit, used by \SI{}{\pu})
% =====================================================================

\section{Adding the Tower Structural Modes to the Reduced Wind-Turbine Model}
\label{sec:simplified-tower-model}

The OpenFAST co-simulation of \S\ref{sec:openfast-excitation} and
\S\ref{sec:tower-ss-resonance} established, in a high-fidelity aeroelastic model,
that the tower side-to-side (SS) mode can be driven by a generator-torque
oscillation. That model is, however, expensive and its electrical$\to$structure
path is partly obstructed by ROSCO (\S\ref{subsec:ss-torque-path}). To study the
interaction \emph{systematically} --- swept over frequency, contrasted against
the fore-aft (FA) mode, and combined with realistic grid events on the full
LEOGO network --- a fast, fully electrical-to-mechanical model is needed. This
section describes such a model: \texttt{WindTurbineTower}, a superset of the
reduced two-mass turbine of \S\ref{sec:frequency-analysis} that adds the first
tower side-to-side and first tower fore-aft bending modes as modal degrees of
freedom.

\subsection{Modal (assumed-mode) representation}
\label{subsec:stm-modal}

Each tower mode is represented as a single-degree-of-freedom modal oscillator,
exactly as ElastoDyn/FAST represents tower bending through a Rayleigh--Ritz mode
shape. The side-to-side displacement \(q_{\mathrm{ss}}\) obeys
\begin{equation}
  \ddot{q}_{\mathrm{ss}}
    + 2\zeta_{\mathrm{ss}}\omega_{\mathrm{ss}}\,\dot{q}_{\mathrm{ss}}
    + \omega_{\mathrm{ss}}^2\,q_{\mathrm{ss}}
  = g_{\mathrm{ss}}\,T_e,
  \qquad \omega_{\mathrm{ss}} = 2\pi f_{\mathrm{ss}},
  \label{eq:stm-ss}
\end{equation}
and the tower-top side-to-side \emph{acceleration} (the analogue of the OpenFAST
output \texttt{YawBrTAyp}) is its second derivative
\begin{equation}
  a_{\mathrm{ss}} = \ddot{q}_{\mathrm{ss}}
    = g_{\mathrm{ss}}\,T_e
      - 2\zeta_{\mathrm{ss}}\omega_{\mathrm{ss}}\,\dot{q}_{\mathrm{ss}}
      - \omega_{\mathrm{ss}}^2\,q_{\mathrm{ss}}.
  \label{eq:stm-ss-accel}
\end{equation}
Here \(T_e = P_e/(\omega_{e,\mathrm{filt}}\,\eta)\) is the same electromagnetic
torque that drives the drivetrain equation \eqref{eq:Te}, so the driving term of
\eqref{eq:stm-ss} carries the grid disturbance directly. The physical mechanism
is the nacelle reaction torque: the air-gap torque \(T_e\) acts about the
low-speed-shaft axis, and by Newton's third law the nacelle feels the equal and
opposite reaction, which rolls it laterally and bends the tower top side to side.
This is the same mechanism exploited by active side-to-side tower-damping
controllers and by ROSCO's Tower Resonance Avoidance logic, both of which
modulate generator torque to influence tower motion.

The fore-aft displacement \(q_{\mathrm{fa}}\) mirrors \eqref{eq:stm-ss} but is
driven by the rotor \emph{thrust} rather than the torque,
\begin{equation}
  \ddot{q}_{\mathrm{fa}}
    + 2\zeta_{\mathrm{fa}}\omega_{\mathrm{fa}}\,\dot{q}_{\mathrm{fa}}
    + \omega_{\mathrm{fa}}^2\,q_{\mathrm{fa}}
  = g_{\mathrm{fa}}\,\hat{F}_T,
  \qquad
  \hat{F}_T = \frac{C_T(\lambda,\beta)\,v^2}
                   {\big(C_T(\lambda,\beta)\,v^2\big)_0},
  \label{eq:stm-fa}
\end{equation}
where \(\hat{F}_T\) is the thrust normalised by its equilibrium value (so a
steady thrust produces no ring), \(\lambda = \omega_m R/v\) is the tip-speed
ratio and \(\beta\) the blade pitch. The two oscillators add four states
\(\{q_{\mathrm{ss}},\dot{q}_{\mathrm{ss}},q_{\mathrm{fa}},\dot{q}_{\mathrm{fa}}\}\)
to the model.

\subsection{Forward coupling: the tower does not react back}
\label{subsec:stm-forward}

The two modal oscillators are driven by the drivetrain/aerodynamic state
(\(T_e\) for SS, \(\hat{F}_T\) for FA) but do \emph{not} feed back onto the
drivetrain or the electrical network: they are pure output DOFs. This
one-directional (forward) coupling is deliberate and has two consequences. First,
it is physically reasonable here --- the tower-top lateral acceleration is orders
of magnitude too small to modify the air-gap torque or the delivered power.
Second, and more usefully for the study, it means that switching a mode on or off
(\(\texttt{ss\_enable}\), \(\texttt{fa\_enable}\)) leaves the electrical solution
\emph{bit-for-bit identical}. Any two runs that differ only in the tower DOFs
therefore share exactly the same grid trajectory, so a comparison isolates the
mechanical response without confounding electrical differences. The same holds
for on/off comparisons of the frequency-support droop (\S\ref{sec:outer-droop}),
which is why the droop trade-off of \S\ref{subsec:nt-tradeoff} is clean.

\subsection{Calibration against the OpenFAST-FMU}
\label{subsec:stm-calibration}

The modal frequencies and damping ratios are taken from, and the coupling gains
calibrated against, the high-fidelity OpenFAST model, so that the reduced model
reproduces the FMU tower response quantitatively rather than only qualitatively.

\paragraph{Frequencies.}
The IEA 15\,MW monopile first tower modes sit at
\(f_{\mathrm{ss}} = \SI{0.234}{\hertz}\) (side-to-side) and
\(f_{\mathrm{fa}} = \SI{0.235}{\hertz}\) (fore-aft); the two are nearly
degenerate in frequency, which makes the \emph{drive path}, not the frequency,
the discriminator between them (\S\ref{sec:shaft-excitation}).

\paragraph{Side-to-side gain \(g_{\mathrm{ss}}\).}
For the oscillator \eqref{eq:stm-ss} driven at its natural frequency, the
steady-state acceleration amplitude follows from the transfer function of
\eqref{eq:stm-ss-accel} at \(\omega=\omega_{\mathrm{ss}}\),
\begin{equation}
  \big|a_{\mathrm{ss}}(\omega_{\mathrm{ss}})\big|
  = \frac{g_{\mathrm{ss}}}{2\zeta_{\mathrm{ss}}}\,\big|\Delta T_e\big|.
  \label{eq:stm-gss}
\end{equation}
The OpenFAST-FMU torque sweep (\S\ref{sec:tower-ss-resonance}) measured a peak
side-to-side acceleration of \(\approx\SI{0.259}{\metre\per\second\squared}\)
(\texttt{YawBrTAyp} at \SI{0.234}{\hertz}) for a \(\pm\SI{10}{\percent}\)
generator-torque modulation (\(|\Delta T_e|\approx\SI{0.10}{\pu}\)) at
\(\zeta_{\mathrm{ss}}\approx\num{0.05}\). Solving \eqref{eq:stm-gss} gives a
first estimate \(g_{\mathrm{ss}}\approx\SI{0.26}{\metre\per\second\squared\per\pu}\),
which is then refined by an equivalent reduced-model torque sweep
(\texttt{sweep\_ss\_reduced.py}) that injects the same \(\pm\SI{10}{\percent}\)
modulation on \(T_e\) and matches the FMU peak, yielding the calibrated default
\begin{equation}
  g_{\mathrm{ss}} = \SI{0.2684}{\metre\per\second\squared\per\pu},
  \qquad
  \zeta_{\mathrm{ss}} = \num{0.05}.
  \label{eq:stm-gss-value}
\end{equation}

\paragraph{Fore-aft gain and effective damping.}
The fore-aft gain \(g_{\mathrm{fa}} = \SI{0.102}{\metre\per\second\squared\per\pu}\)
and effective damping \(\zeta_{\mathrm{fa}} = \num{0.0125}\) are calibrated the
same way, from the FMU fore-aft ring-down and forced-response at
\(v = \SI{11}{\metre\per\second}\). The damping is termed \emph{effective}
because the strong aerodynamic damping of the fore-aft mode --- when the tower
top moves downwind the rotor sees less wind, the thrust drops and opposes the
motion --- is folded into a single modal \(\zeta_{\mathrm{fa}}\) rather than
resolved explicitly. Notably \(\zeta_{\mathrm{fa}}\) here is \emph{small}: the
reason the fore-aft mode is not the resonance-prone one is not heavy damping but
the weak thrust drive path (\S\ref{sec:shaft-excitation}).

Table~\ref{tab:stm-params} collects the calibrated modal parameters. With them
the reduced model reproduces the FMU side-to-side resonance peak to within a few
percent (\S\ref{sec:new-tests}) while running in real time against the full
LEOGO network.

\begin{table}[t]
  \centering
  \caption{Calibrated tower modal parameters of \texttt{WindTurbineTower},
  taken from / matched to the OpenFAST-FMU IEA 15\,MW monopile model.}
  \label{tab:stm-params}
  \begin{tabular}{lccc}
    \toprule
    Mode & \(f\) [\si{\hertz}] & \(\zeta\) [--] &
      \(g\) [\si{\metre\per\second\squared\per\pu}] \\
    \midrule
    Side-to-side (torque-driven) & 0.234 & 0.050 & 0.2684 \\
    Fore-aft (thrust-driven)     & 0.235 & 0.0125 & 0.102 \\
    \bottomrule
  \end{tabular}
\end{table}

\subsection{Ring-down time constant}
\label{subsec:stm-tau}

A consequence of the light side-to-side damping worth stating explicitly, since
it recurs in the scenarios of \S\ref{sec:new-tests}, is the modal decay time
constant
\begin{equation}
  \tau_{\mathrm{ss}} = \frac{1}{\zeta_{\mathrm{ss}}\,\omega_{\mathrm{ss}}}
    = \frac{1}{0.05\cdot 2\pi\cdot 0.234}
    \approx \SI{13.6}{\second}.
  \label{eq:stm-tau}
\end{equation}
A broadband ``kick'' (a load step) therefore excites a free side-to-side
ring-down that decays with \(\tau_{\mathrm{ss}}\approx\SI{13.6}{\second}\),
whereas a sustained sinusoidal drive at \(f_{\mathrm{ss}}\) builds the mode up
over \(\sim 3\text{--}4\,\tau_{\mathrm{ss}}\approx\SI{45}{\second}\) to a large
steady amplitude (\S\ref{subsec:nt-resonance}). The distinction between the two
--- ring-\emph{down} versus build-\emph{up} --- is the clearest experimental
signature of resonant versus broadband excitation.

\subsection*{Sources}
\label{subsec:stm-sources}

\begin{itemize}
  \item Modal (assumed-mode / Rayleigh--Ritz) tower-bending DOFs, as used by
        ElastoDyn/OpenFAST: Jonkman (2007), \emph{Dynamics Modeling and Loads
        Analysis of an Offshore Floating Wind Turbine}, NREL/TP-500-41958,
        \href{https://www.nrel.gov/docs/fy08osti/41958.pdf}{nrel.gov/docs/fy08osti/41958.pdf};
        Jonkman \& Buhl (2005), \emph{FAST User's Guide}, NREL/EL-500-38230,
        \href{https://www.nrel.gov/docs/fy06osti/38230.pdf}{nrel.gov/docs/fy06osti/38230.pdf};
        OpenFAST documentation, \href{https://openfast.readthedocs.io/}{openfast.readthedocs.io}.
  \item IEA \SI{15}{\mega\watt} reference turbine (modal frequencies, rotor and
        tower data): Gaertner et al.\ (2020), \emph{Definition of the IEA
        15-Megawatt Offshore Reference Wind Turbine}, NREL/TP-5000-75698,
        \href{https://www.nrel.gov/docs/fy20osti/75698.pdf}{nrel.gov/docs/fy20osti/75698.pdf};
        model repository
        \href{https://github.com/IEAWindTask37/IEA-15-240-RWT}{github.com/IEAWindTask37/IEA-15-240-RWT}.
  \item Fore-aft aerodynamic tower damping (the ``aerodynamic damper''):
        K\"uhn (2001), \emph{Dynamics and Design Optimisation of Offshore Wind
        Energy Conversion Systems}, PhD thesis, TU Delft; Salzmann \& van der
        Tempel (2005), \emph{Aerodynamic damping in the design of support
        structures for offshore wind turbines}, Copenhagen Offshore Wind
        Conference.
\end{itemize}

% =====================================================================
%  Seksjon: Kontroller-blokkering -- hvorfor ROSCO hindrer at en
%           nett-/elektrisk forstyrrelse eksiterer tarnmodene (side-to-side
%           OG fore-aft) i den hoy-fidelitets OpenFAST+ROSCO-turbinen.
%           Samler bevis-testene og forklarer blokkeringen mode-for-mode.
%  Kilder:  docs/openfast_excitation.tex        (detaljert SS-behandling)
%           docs/STATUS_OG_VEIVALG.md, kap. 2-3  (testtabell + rotarsak)
%           test1002/ControlData/ROSCO.IEA15MW.IN (VSContrl, OL_Mode, TRA_Mode,
%                                                   F_LPFCornerFreq)
%  Preamble: amsmath, amssymb, graphicx, booktabs, siunitx, hyperref
% =====================================================================

\section{Controller Blocking: Why ROSCO Prevents Electrical Excitation of the Tower Modes}
\label{sec:rosco-blocking}

This project studies electro-mechanical interactions in offshore wind turbines
coupled to a grid --- here an IEA \SI{15}{\mega\watt} turbine connected to the
LEOGO islanded platform network. One such interaction is whether an
electrical/grid disturbance can reach the turbine's structural tower modes. In
the high-fidelity OpenFAST turbine that particular path fails: neither the
side-to-side (SS) nor the fore-aft (FA) tower mode can be excited through any
electrical input, because the ROSCO controller owns the one channel that would
carry the disturbance into the structure. This section gathers
the tests that established this negative result --- the tests that ``proved'' the
blocking --- for \emph{both} tower modes, and explains mode-by-mode why the
electrical path is closed. It is a pivot point of the methodology: it is precisely
this blocking that motivates the ROSCO open-loop torque injection that bypasses it
(\S\ref{subsec:se-openfast}) and the reduced \texttt{WindTurbineTower} whose
electrical-to-mechanical path is open by construction
(\S\ref{sec:simplified-tower-model}). The detailed SS experiment log lives in
\S\ref{sec:openfast-excitation}; here the focus is the evidence and the causal
argument.

\subsection{The two candidate drive paths}
\label{subsec:rb-paths}

For a grid disturbance to reach a tower mode it must modulate one of two
mechanical loads: the \emph{generator torque}, whose nacelle reaction moment rolls
the nacelle laterally and bends the tower side to side; or the \emph{rotor
thrust}, which bends the tower fore and aft. The two first tower modes are nearly
degenerate in frequency
(\(f_{\mathrm{ss}}\approx f_{\mathrm{fa}}\approx\SI{0.234}{\hertz}\)), so what
decides whether an electrical disturbance can reach each of them is not the
frequency but whether the disturbance can act on torque (for SS) or on thrust
(for FA). Table~\ref{tab:rb-paths} states the mapping that the tests below probe.

\begin{table}[t]
  \centering
  \caption{The two electrical-to-structure drive paths and the controller that
  guards each in the OpenFAST-FMU (\texttt{VSContrl}=5).}
  \label{tab:rb-paths}
  \begin{tabular}{lll}
    \toprule
    Tower mode & Mechanical drive & Grid access in the FMU \\
    \midrule
    Side-to-side & generator torque \(T_e\) & only via ROSCO's torque command \\
    Fore-aft     & rotor thrust \(F_T\)      & no electrical thrust channel at all \\
    \bottomrule
  \end{tabular}
\end{table}

That the torque-to-tower coupling itself \emph{exists} in the model is not in
doubt: ROSCO even ships a Tower Resonance Avoidance mode (\texttt{TRA\_Mode}) that
deliberately modulates \emph{generator torque} to steer the rotor away from a
tower frequency. The coupling is real; the question the tests answer is whether an
\emph{external, electrically imposed} torque or thrust can use it.

\subsection{Tests that closed the torque path (side-to-side)}
\label{subsec:rb-torque}

Three tests, each a baseline-versus-modulated comparison, established that the
generator-torque channel is closed to external commands.

\begin{enumerate}
  \item \textbf{Generator-torque command} (\texttt{GenSpdOrTrq}). A commanded
        torque modulation at the SS frequency left the simulated generator torque
        \emph{bit-identical} to the unmodulated baseline --- the command has no
        effect whatsoever.
  \item \textbf{Electrical-power command} (\texttt{ElecPwrCom},
        \(\pm\SI{20}{\percent}\) at \SI{0.236}{\hertz}). Ignored identically in
        Region~2 (\SI{8}{\metre\per\second}) and Region~3: the \texttt{GenTq}
        ripple and the SS acceleration envelope were indistinguishable from
        baseline (\S\ref{subsec:of-electrical}, Table therein).
  \item \textbf{Opening the external-torque path} (\texttt{VSContrl}=4). The one
        configuration that would let an external torque become the actual
        generator torque \emph{crashes} the FMU at the first \texttt{doStep}
        (access-violation in the FMI call); the binary lacks the
        Simulink/Labview pointer hookup and it is not fixable from configuration
        (\S\ref{subsec:of-vscontrl4}).
\end{enumerate}

The cause is architectural. With \(\texttt{VSContrl}=5\) ROSCO computes the
generator torque \emph{internally} from the measured generator speed --- in
Region~2 the optimal law
\begin{equation}
  T_{\mathrm{gen}} = K_{\mathrm{opt}}\,\omega_g^2,
  \label{eq:rb-komega2}
\end{equation}
and does not read the external \texttt{GenSpdOrTrq}/\texttt{ElecPwrCom} ports. The
small torque ripple seen in the runs is ROSCO's own speed-reaction to the start-up
transient, not a response to the command. Table~\ref{tab:rb-evidence} collects the
evidence. The verdict for the torque path, and hence for the SS mode, is
unambiguous: \emph{it is closed}.

\begin{table}[t]
  \centering
  \caption{The controller-blocking evidence (OpenFAST-FMU, \texttt{VSContrl}=5).
  Every electrical route into the generator torque is either ignored or crashes.}
  \label{tab:rb-evidence}
  \begin{tabular}{lll}
    \toprule
    Test & Result & Root cause \\
    \midrule
    Torque cmd.\ \texttt{GenSpdOrTrq}      & bit-identical \texttt{GenTq}        & ROSCO owns torque (\ref{eq:rb-komega2}) \\
    Power cmd.\ \texttt{ElecPwrCom} \(\pm20\%\) & ignored, Region~2 \emph{and} 3 & same --- torque set internally \\
    \texttt{VSContrl}=4 (ext.\ torque)     & crash at first \texttt{doStep}      & missing Simulink pointer; needs rebuild \\
    ROSCO filters hiding SS?               & no (see \S\ref{subsec:rb-filters})  & filters shape response, not access \\
    \bottomrule
  \end{tabular}
\end{table}

\subsection{Ruling out the ROSCO filters as the cause}
\label{subsec:rb-filters}

A natural suspicion is that the blocking is really a \emph{filtering} effect ---
that ROSCO's low-pass and notch filters (\texttt{F\_LPFCornerFreq} and the tower
notch used by the SS-damping / resonance-avoidance logic) are actively cancelling
the SS excitation before it reaches the structure. This was tested and rejected.
Those filters act only on the \emph{speed signal that ROSCO measures}, shaping
\emph{how ROSCO responds}; they are not the reason the network fails to reach the
SS mode. The disturbance is stopped one step earlier, at the torque command
itself: because \(\texttt{VSContrl}=5\) makes the torque a pure function of
measured speed \eqref{eq:rb-komega2}, the external command never enters the torque
in the first place, so there is nothing for a notch filter to remove. The blocking
is architectural, not spectral --- an important distinction, because it means no
filter re-tuning could open the path; only a different converter interface could.

\subsection{The fore-aft mode is doubly unreachable}
\label{subsec:rb-foreaft}

The fore-aft mode is blocked even more firmly than the side-to-side mode, for a
different reason. FA is driven by rotor \emph{thrust}, and the FMU exposes no
electrical input that acts on thrust at all --- there is simply no
grid-to-thrust channel to command. The only conceivable route is the indirect one
already written in \eqref{eq:se-fa-path},
\[
  \text{grid}\to P_e\to T_e\to \omega_m\to \lambda\to C_T\to F_T
  \to q_{\mathrm{fa}},
\]
and every gate on that route is shut: its first link is the very torque channel
that ROSCO owns (\S\ref{subsec:rb-torque}), so \(T_e\) barely moves; and even if it
did, the huge rotor inertia would keep \(\omega_m\), and therefore \(\lambda\),
\(C_T\) and the thrust, nearly constant at fixed wind. The FA mode is thus blocked
twice over --- no direct electrical thrust channel, and the indirect
torque\(\to\)speed\(\to\)thrust channel throttled by both ROSCO and the rotor
inertia. Neither tower mode, then, can be electrically excited in the FMU.

This is also why the FA/SS distinction survives into the reduced model, where the
torque path \emph{is} open: there SS is driven and FA is not, by a factor of
\(\sim\!30\) at resonance (\S\ref{subsec:se-selectivity},
\S\ref{subsec:nt-resonance}). The controller blocking in the FMU and the drive-path
selectivity in the reduced model are two views of the same fact --- the grid can
only ever reach the tower through torque, and only the side-to-side mode listens to
torque.

\subsection{The core: full-converter decoupling}
\label{subsec:rb-core}

Putting the pieces together, the OpenFAST-FMU behaves as a \emph{full-converter
decoupling with no return path}. With \(\texttt{VSContrl}=5\) the OpenFAST
generator torque is entirely ROSCO's; the network disturbance reaches only the
co-simulation \emph{wrapper} generator mass --- which is why it \emph{does} excite
the \SI{3.49}{\hertz} drivetrain torsional mode of \S\ref{sec:forcedresponse} ---
but never OpenFAST's own structural solver, and hence never the tower modes. The
electrical and structural sides of the high-fidelity model are separated by the
controller, and there is no path back.

Two design choices follow directly, and both are pursued elsewhere in this report.
To prove that the OpenFAST \emph{tower} really would amplify a torque oscillation
if one could be imposed, the torque is injected through ROSCO's own open-loop
control table, bypassing the \(K\omega^2\) law entirely
(\S\ref{subsec:se-openfast}); that experiment recovers the sharp SS resonance and
so isolates the structural selectivity from the controller blocking. And to obtain
a genuinely \emph{closed} electrical drive of the tower --- a real grid
disturbance producing the torque --- the study moves to the reduced
\texttt{WindTurbineTower} (\S\ref{sec:simplified-tower-model}), where
\(T_e = P_e/(\omega_{e,\mathrm{filt}}\eta)\) enters the modal equation explicitly
and no controller intercepts it. The blocking result of this section is thus not a
dead end but the hinge of the whole methodology.

\subsection*{Sources}
\label{subsec:rb-sources}

\begin{itemize}
  \item ROSCO controller (\texttt{VSContrl} torque law, open-loop control,
        \texttt{TRA\_Mode} tower resonance avoidance, \texttt{F\_LPFCornerFreq}):
        Abbas et al.\ (2022), \emph{A reference open-source controller for fixed
        and floating offshore wind turbines}, Wind Energy Science 7, 53--73,
        \href{https://doi.org/10.5194/wes-7-53-2022}{doi.org/10.5194/wes-7-53-2022};
        documentation \href{https://rosco.readthedocs.io/}{rosco.readthedocs.io};
        source \href{https://github.com/NREL/ROSCO}{github.com/NREL/ROSCO}.
  \item OpenFAST / ServoDyn (\texttt{VSContrl} options, FMU interface):
        \href{https://openfast.readthedocs.io/}{openfast.readthedocs.io}.
\end{itemize}

% =====================================================================
%  Seksjon: Skaft-/generatormoment-eksitasjon av tårnmodene i begge
%           modeller (forenklet WindTurbineTower + OpenFAST-FMU), inkl.
%           manipulering av ROSCO open-loop-fila, og SS-vs-fore-aft-
%           selektiviteten.
%  Kilder:  src/tops_openfast/dyn_models/windturbine_tower.py
%           casestudies/dyn_sim/sweep_resonance.py (FMU, OL-table)
%           casestudies/dyn_sim/_contrast_ss_fa.py
%           test1002/ControlData/ROSCO.IEA15MW.IN  (OL_Mode, Ind_GenTq)
%  Preamble: amsmath, amssymb, graphicx, booktabs, siunitx, hyperref
% =====================================================================

\section{Shaft (Generator-Torque) Excitation of the Tower Modes}
\label{sec:shaft-excitation}

This project studies electro-mechanical interactions in offshore wind turbines
coupled to a grid --- here one connected to the LEOGO islanded platform network.
One important such interaction is whether an \emph{electrical} disturbance can
excite a \emph{mechanical} tower mode. The physical link is the generator/air-gap
electromagnetic torque acting through the shaft: fluctuations in the delivered
power \(P_e\) become fluctuations in \(T_e = P_e/(\omega_{e,\mathrm{filt}}\eta)\),
whose nacelle reaction torque bends the tower side to side. This section collects
how that shaft-torque excitation is realised in the two models --- the reduced
\texttt{WindTurbineTower} of \S\ref{sec:simplified-tower-model} and the
OpenFAST-FMU of \S\ref{sec:openfast-excitation} --- and shows that the torque
path selectively reaches the side-to-side (SS) mode but not the fore-aft (FA)
mode.

\subsection{Reduced model: an open electrical-to-mechanical path}
\label{subsec:se-reduced}

In \texttt{WindTurbineTower} the side-to-side oscillator \eqref{eq:stm-ss} is
driven directly by \(T_e\), and \(T_e\) is computed from the power the converter
actually delivers to the grid. The path
\begin{equation}
  \text{grid disturbance}\;\longrightarrow\;P_e
  \;\longrightarrow\;T_e
  \;\longrightarrow\;\text{nacelle reaction moment}
  \;\longrightarrow\;q_{\mathrm{ss}}
  \label{eq:se-path}
\end{equation}
is therefore \emph{open by construction}: no controller intercepts it, unlike in
the FMU. A sustained sinusoidal grid load at the modal frequency thus feeds the
SS mode continuously and drives it into resonance, and a swept load traces out
its resonance curve. Both are demonstrated quantitatively in
\S\ref{sec:new-tests}; here the point is only that the reduced model provides the
genuinely closed electrical\(\to\)mechanical demonstration that the FMU cannot,
because in the FMU ROSCO owns the torque.

\subsection{OpenFAST model: injecting torque through the ROSCO open-loop file}
\label{subsec:se-openfast}

In the OpenFAST build used here ROSCO (\texttt{VSContrl}=5) computes the
generator torque internally from the measured generator speed and rejects every
external torque or power command passed across the FMU interface
(\texttt{GenSpdOrTrq}, \texttt{ElecPwrCom}); this is the negative result
documented in \S\ref{subsec:of-electrical}. Consequently the network cannot reach
the OpenFAST \emph{structure} through the normal electrical inputs, and the only
channel that does reach the structural solver is ROSCO's own \emph{open-loop
control} table.

\paragraph{What was changed in the ROSCO input.}
ROSCO reads its input files at runtime, so the torque channel was opened without
rebuilding the FMU, by editing \texttt{test1002/ControlData/ROSCO.IEA15MW.IN}:
\begin{itemize}
  \item \texttt{OL\_Mode}: \(0\to1\) --- enable open-loop control versus time.
  \item \texttt{Ind\_GenTq}: \(0\to2\) --- take the generator torque from column~2
        (in \si{\newton\metre}) of the open-loop table, with time in column~1
        (\texttt{Ind\_Breakpoint}=1).
  \item A new table file \texttt{OLInput\_ROSCO.dat} was written with the
        prescribed torque time series.
\end{itemize}
The prescribed torque is a sinusoidal modulation about the Region-2 operating
torque \(T_0\),
\begin{equation}
  T_{\mathrm{gen}}(t) = T_0\,\big[\,1 + a\,
    \sin\!\big(2\pi f\,(t-t_{\mathrm{on}})\big)\,\big]
    \mathbb{1}[t\ge t_{\mathrm{on}}],\qquad a = 0.10,
  \label{eq:se-ol-forcing}
\end{equation}
i.e. a \(\pm\SI{10}{\percent}\) generator-torque oscillation. This overrides
ROSCO's torque regulation entirely, so the runs are kept short and the rotor
speed monitored for drift, with \(T_0\) matched to the aerodynamic torque at the
operating wind speed so that the net drift over the window is negligible.

\paragraph{Did it work, and what did it show.}
It worked, and it is the mechanism behind the frequency sweep of
\S\ref{sec:tower-ss-resonance}. Injecting \eqref{eq:se-ol-forcing} at the modal
frequency produced a tower-top side-to-side acceleration (\texttt{YawBrTAyp})
that \emph{builds up} to a sustained forced amplitude, rather than the free decay
seen in all the ignored \texttt{GenSpdOrTrq}/\texttt{ElecPwrCom} attempts ---
direct confirmation that the torque now reaches the structure. Sweeping \(f\)
gave a sharp resonance peaking in the SS band
(\SIrange{0.234}{0.245}{\hertz}), with a peak side-to-side acceleration of
\(\approx\SI{0.26}{\metre\per\second\squared}\) and roughly \(60\times\)
amplification over the off-resonance floor (Table~\ref{tab:ss-sweep}). The
important caveat, stated already in \S\ref{subsec:ss-torque-path}, is that the
torque here is \emph{imposed} through the open-loop table rather than produced by
a closed grid\(\to\)torque feedback: the experiment proves the \emph{structural}
frequency selectivity of the OpenFAST tower, while the reduced model
(\S\ref{subsec:se-reduced}) supplies the genuinely closed electrical drive. The
two are complementary --- the FMU shows the real tower amplifies a torque
oscillation, and the reduced model shows the grid can supply that oscillation.

\subsection{Why side-to-side and not fore-aft}
\label{subsec:se-selectivity}

The two first tower modes are nearly degenerate in frequency
(\(f_{\mathrm{ss}}\approx f_{\mathrm{fa}}\approx\SI{0.234}{\hertz}\)), so what
separates them under an electrical disturbance is not frequency but the
\emph{drive path}. The side-to-side mode is driven by the generator torque, which
carries the grid disturbance directly (\eqref{eq:se-path}). The fore-aft mode is
driven by the rotor thrust \(F_T\), which a grid disturbance can reach only
indirectly, through the drivetrain speed,
\begin{equation}
  \text{grid}\to P_e\to T_e\to \omega_m\to \lambda\to C_T\to F_T
  \to q_{\mathrm{fa}},
  \label{eq:se-fa-path}
\end{equation}
i.e. a second-order path gated by the very large rotor inertia. At a fixed wind
speed the thrust barely moves when the grid is perturbed, so the fore-aft
response to an electrical disturbance is much smaller than the side-to-side
response. This asymmetry is exactly why the SS mode --- and not the frequency-
adjacent FA mode --- is the electrically resonance-prone one.

The reduced model makes this quantitative. Because the two modes are switched
independently and share the same electrical solution
(\S\ref{subsec:stm-forward}), a single co-simulation exposes both responses to
the identical disturbance. Under a sustained load oscillation at
\SI{0.234}{\hertz} the side-to-side lock-in acceleration exceeds the fore-aft one
by a factor of \(\sim 30\) (\S\ref{subsec:nt-resonance}), and under a discrete
grid event the ratio of peak accelerations is of the same order
(\S\ref{subsec:nt-gttrip}). The tower acts as a mechanical band-pass filter on
the electrical disturbance, and the generator-torque path selects the
side-to-side mode as its output. A lock-in / contrast analysis of the two
accelerations at the forcing frequency is provided by
\texttt{casestudies/dyn\_sim/\_contrast\_ss\_fa.py}.

\subsection*{Sources}
\label{subsec:se-sources}

\begin{itemize}
  \item ROSCO open-loop control (\texttt{OL\_Mode}, \texttt{Ind\_GenTq},
        \texttt{OLInput\_ROSCO.dat}) and the internal \(K\omega^2\) torque law:
        Abbas et al.\ (2022), \emph{A reference open-source controller for fixed
        and floating offshore wind turbines}, Wind Energy Science 7, 53--73,
        \href{https://doi.org/10.5194/wes-7-53-2022}{doi.org/10.5194/wes-7-53-2022};
        documentation \href{https://rosco.readthedocs.io/}{rosco.readthedocs.io}
        (see the \emph{Open Loop Input} page); source
        \href{https://github.com/NREL/ROSCO}{github.com/NREL/ROSCO}.
  \item OpenFAST (FMU export, \texttt{YawBrTAyp} tower-top acceleration output):
        \href{https://openfast.readthedocs.io/}{openfast.readthedocs.io}.
\end{itemize}

% =====================================================================
%  Seksjon: Den ytre frekvensstøtte-droopen i vindturbinmodellen.
%  Kilde:   src/tops_openfast/dyn_models/windturbine_tower.py
%             _droop_command(), P_ref_components(), set_grid_frequency_hz()
%  Driver:  casestudies/dyn_sim/test_WT_LEOGO_tower_sim.py  (feed_grid_frequency)
%  Preamble: amsmath, amssymb, graphicx, booktabs, siunitx
%           + \DeclareSIUnit{\pu}{pu}   (per-unit, used by \SI{}{\pu})
% =====================================================================

\section{The Outer Wind-Turbine Frequency-Support Droop}
\label{sec:outer-droop}

To place the tower-mode interaction in an operationally realistic setting, the
reduced turbine of \S\ref{sec:simplified-tower-model} carries an optional
\emph{frequency-support droop}: on an islanded platform grid such as LEOGO the
wind turbine is expected to help arrest frequency excursions after a disturbance,
the same service the gas-turbine governors provide. This section defines that
outer droop law; whether it is actually needed, given that the converter already
contains a droop, is the subject of \S\ref{sec:outer-droop-necessity}.

\subsection{Control law}
\label{subsec:od-law}

The droop is implemented in \texttt{WindTurbineTower.\_droop\_command}. The
turbine is first \emph{de-loaded}: instead of tracking the maximum-power-tracking
(MPT) point \(P_{\mathrm{avail}}\), it holds a reduced base reference
\begin{equation}
  P_{\mathrm{base}} = \operatorname{clip}\!\big(
    P_{\mathrm{avail}} - P_{\mathrm{hr}},\;0,\;P_{\mathrm{avail}}\big),
  \label{eq:od-base}
\end{equation}
keeping a headroom reserve \(P_{\mathrm{hr}}\) (parameter
\(\texttt{headroom\_pu}\), on the UIC base) available for upward support. On top
of the base it adds a proportional response to the measured grid-frequency
deviation,
\begin{equation}
  \Delta P = K_{\mathrm{droop}}\,\big(f_{\mathrm{nom}} - f_{\mathrm{grid}}\big),
  \qquad
  P_{\mathrm{ref}} = \operatorname{clip}\!\big(
    P_{\mathrm{base}} + \Delta P,\;0,\;P_{\mathrm{avail}}\big),
  \label{eq:od-delta}
\end{equation}
so that a frequency dip (\(f_{\mathrm{grid}} < f_{\mathrm{nom}}\)) raises the
power reference and injects extra active power into the grid, up to the headroom
limit. The gain \(K_{\mathrm{droop}}\) is \(\texttt{k\_droop\_pu\_per\_hz}\) and
\(f_{\mathrm{nom}} = \SI{50}{\hertz}\). The droop is toggled with
\(\texttt{droop\_enable}\); when it is off, \(\Delta P = 0\) and the turbine
simply holds the de-loaded base \eqref{eq:od-base}.

An important detail for the experiments is that the headroom \eqref{eq:od-base}
is applied \emph{whether or not} the droop is enabled. Both the droop-off and
droop-on runs therefore share the same de-loaded operating point, and their
difference isolates the droop \emph{response} \eqref{eq:od-delta} alone, with no
confounding change in the base loading.

\subsection{Measured grid frequency}
\label{subsec:od-freq}

The frequency the droop reacts to is supplied from the network side of the
co-simulation through \texttt{set\_grid\_frequency\_hz}. In the LEOGO driver the
platform is dominated by the three coherent gas-turbine synchronous machines on
the main bus, so the frequency is taken as the centre-of-inertia (mean
generator-speed) frequency
\begin{equation}
  f_{\mathrm{grid}} = f_n\big(1 + \overline{\Delta\omega}_{\mathrm{gen}}\big),
  \label{eq:od-fgrid}
\end{equation}
with \(f_n = \SI{50}{\hertz}\) and \(\overline{\Delta\omega}_{\mathrm{gen}}\) the
mean per-unit generator-speed deviation. Because the generator speed is a state,
\(f_{\mathrm{grid}}\) is available before the algebraic solve at each integration
stage and is fed to the turbine, which reads it when it forms \(P_{\mathrm{ref}}\)
(\texttt{feed\_grid\_frequency} in the driver). This closes an outer,
network-measured loop around the turbine, distinct from the converter's own local
droop (\S\ref{sec:outer-droop-necessity}).

\subsection{Parameters and defaults}
\label{subsec:od-params}

Table~\ref{tab:od-params} lists the droop parameters. In the headline scenario of
\S\ref{subsec:nt-tradeoff} the droop uses \(K_{\mathrm{droop}} =
\SI{0.75}{\pu\per\hertz}\) with a \SI{10}{\percent} headroom reserve; these are
the values for which the frequency-support benefit and the tower-loading penalty
are reported.

\begin{table}[t]
  \centering
  \caption{Outer frequency-support droop parameters
  (\texttt{WindTurbineTower}).}
  \label{tab:od-params}
  \begin{tabular}{lll}
    \toprule
    Parameter & Symbol & Value \\
    \midrule
    Nominal frequency        & \(f_{\mathrm{nom}}\)      & \SI{50}{\hertz} \\
    Droop gain               & \(K_{\mathrm{droop}}\)    & \SI{0.75}{\pu\per\hertz} \\
    Headroom reserve         & \(P_{\mathrm{hr}}\)       & \SI{0.10}{\pu} (UIC base) \\
    Enable flag              & \(\texttt{droop\_enable}\) & \(\{0,1\}\) \\
    \bottomrule
  \end{tabular}
\end{table}

% =====================================================================
%  Seksjon: Trenger vi den ytre droopen? To-droop-arkitekturen
%           (UIC-ens innebygde droop vs. den ytre WT-droopen).
%  Kilder:  src/tops_openfast/dyn_models/UIC.py       (perfect_tracking)
%           src/tops_openfast/dyn_models/windturbine_tower.py (_droop_command)
%  Bakgrunn: docs/STATUS_OG_VEIVALG.md, kap. 5
%  Preamble: amsmath, amssymb, graphicx, booktabs, siunitx
% =====================================================================

\section{Do We Need the Outer Droop? The Two-Droop Architecture}
\label{sec:outer-droop-necessity}

The frequency-support droop of \S\ref{sec:outer-droop} is not the only droop in
the model. The converter interface itself --- the Unified Integral Control (UIC)
grid-forming model --- \emph{is} a droop regulator by design, and it is active in
every run. It is therefore fair to ask whether the outer wind-turbine droop is
needed at all, or whether it merely duplicates what the converter already does.
This section separates the two mechanisms, explains what each contributes, and
states when the outer droop earns its place.

\subsection{The inner droop: UIC grid-forming control}
\label{subsec:odn-inner}

The UIC (\texttt{UIC\_sig}) implements the control law
\begin{equation}
  \dot{v}_i = j\,\omega_n K_i\,\varepsilon_s
            + \omega_n K_v\,\varepsilon_v
            + j\,v_i\,x_{\mathrm{filter}}\,\texttt{perfect\_tracking},
  \label{eq:odn-uic}
\end{equation}
where \(\varepsilon_s\) carries the active/reactive-power error and
\(\varepsilon_v\) the voltage error. The \(j\)-rotation turns a power error into a
frequency/angle adjustment (a P--f droop) and a voltage error into a magnitude
adjustment (a Q--V droop). This is an \emph{inherent, local, grid-forming} droop:
it gives the converter a synchronising, inertia-like response, and it is the very
reason a grid disturbance reaches the turbine at all rather than seeing a stiff
current source. The \(\texttt{perfect\_tracking}\) term \eqref{eq:odn-uic} can
cancel the droop and make the converter a PLL-free perfect tracker; with
\(\texttt{perfect\_tracking}=0\), as used throughout, the droop is \emph{on}.

Crucially, this inner droop is what makes the whole electro-mechanical study
possible: it is because the converter is droop-coupled, not stiff, that the
network disturbance propagates through \(P_e\) into \(T_e\) and hence into the
tower (\S\ref{sec:shaft-excitation}). Disabling it would decouple the turbine
from the grid and remove the interaction under study, which is why it is kept on
in every run.

\subsection{The outer droop: an explicit, de-loaded support service}
\label{subsec:odn-outer}

The outer droop \eqref{eq:od-delta} is a different object. It is not a
grid-forming primitive but an \emph{explicit dispatchable service} layered on the
power reference: it de-loads the turbine to hold a headroom reserve
\eqref{eq:od-base} and then releases that reserve in proportion to a measured
frequency deviation. The inner UIC droop does neither --- it responds to the
instantaneous power/voltage error at the terminal, but it does not curtail the
turbine to create a reserve, and it does not target a slow bus-frequency
excursion with a scheduled amount of extra energy. In power-system terms the
inner droop supplies the fast synchronising/inertial response, while the outer
droop supplies the slower primary frequency reserve.

\subsection{So: is it needed?}
\label{subsec:odn-verdict}

The honest answer is \emph{it depends on the question being asked}, and it is
worth stating plainly because it shaped the experimental design.

\begin{itemize}
  \item \textbf{For the electro-mechanical coupling to exist: no.} The grid
        reaches the tower through the inner UIC droop regardless of the outer
        loop. All the shaft-excitation and resonance results
        (\S\ref{sec:shaft-excitation}, \S\ref{subsec:nt-resonance}) are run with
        \(\texttt{droop\_enable}=0\); removing the outer droop does not change
        them. The torsional forced-response results of \S\ref{sec:forcedresponse}
        were likewise obtained with the outer droop off.
  \item \textbf{For demonstrating a realistic frequency-support trade-off: yes.}
        The value the outer droop adds is precisely that it lets the turbine
        \emph{act} on a frequency event with a controllable, de-loaded reserve,
        which is what exposes the central operational tension of this thesis: the
        same active-power injection that helps the grid frequency also drives the
        tower side-to-side mode harder. That trade-off (\S\ref{subsec:nt-tradeoff})
        cannot be posed without an explicit, toggleable support law, and the
        inner UIC droop --- always on, not de-loaded, not dispatchable --- does
        not provide it.
\end{itemize}

Two caveats qualify the outer droop and are carried into the results. First, it
reacts to a \emph{global} centre-of-inertia frequency \eqref{eq:od-fgrid}
computed in the driver from the generator speeds, rather than to a locally
measured bus frequency; this is a mild idealisation appropriate to a small,
coherent islanded grid but would need a local frequency estimator on a larger
network. Second, because the inner UIC droop is always active, the outer loop
sits \emph{on top of} an already frequency-responsive converter; the two are
complementary layers (fast local synchronising response plus slow dispatchable
reserve), not two copies of the same function. The recommendation, consistent
with the code as run, is therefore to keep the inner UIC droop on unconditionally
and to enable the outer droop only when a de-loaded primary-reserve service is
part of the scenario being studied.

% =====================================================================
%  Seksjon: Nye co-simuleringsscenarioer på den reduserte modellen mot
%           hele LEOGO-nettet (driftshendelse, frekvensstøtte-avveining,
%           resonant prosesslast, resonanskurve).
%  Driver:  casestudies/dyn_sim/test_WT_LEOGO_tower_sim.py
%  Analyse: casestudies/dyn_sim/_analyze_gt_trip.py
%           casestudies/dyn_sim/_analyze_freq_support.py
%           casestudies/dyn_sim/_analyze_resonance.py
%           casestudies/dyn_sim/sweep_ss_resonance.py
%  Data:    results/tower_test/{gt_trip,gt_trip_support,resonance_0234,
%                               ss_resonance_curve}.{csv,png}
%  Preamble: amsmath, amssymb, graphicx, booktabs, siunitx
%           + \DeclareSIUnit{\pu}{pu}   (per-unit, used by \SI{}{\pu})
% =====================================================================

\section{New Co-Simulation Scenarios on the LEOGO Network}
\label{sec:new-tests}

With the tower modes calibrated (\S\ref{sec:simplified-tower-model}) and the
shaft-torque drive path established (\S\ref{sec:shaft-excitation}), the reduced
\texttt{WindTurbineTower} is exercised against the full LEOGO microgrid in a
common test harness (\texttt{test\_WT\_LEOGO\_tower\_sim.py}). The turbine
replaces the original LEOGO WTG converter at \texttt{Busbar WTG1 LV} through a
UIC interface, and the three gas-turbine generators, their governors and AVRs,
and the platform loads are simulated as usual with a fixed communication step of
\SI{0.01}{\second}. The harness applies disturbances at the main gas-turbine bus
--- a smooth (ramped) load step, a sustained sinusoidal load modulation, or a
short circuit --- and logs the grid frequency, the turbine electrical variables,
and both tower accelerations. Four scenarios are reported, each targeting a
different aspect of the electro-mechanical interaction.

\subsection{A discrete grid event: the tower as a mechanical band-pass filter}
\label{subsec:nt-gttrip}

The first scenario is a realistic single disturbance: a step change in platform
load, representative of a large consumer switching or a generator trip, applied
as a smooth shunt load at the main bus. The step is broadband --- most of its
spectral energy lies near DC --- yet the tower side-to-side response is a clean
decaying oscillation at \(f_{\mathrm{ss}}\approx\SI{0.234}{\hertz}\), the modal
frequency, not at the (much slower) frequency of the grid transient itself. The
step acts as an impulsive ``kick'' that excites the free side-to-side ring-down
of \eqref{eq:stm-tau}, which decays with \(\tau_{\mathrm{ss}}\approx
\SI{13.6}{\second}\). The fore-aft mode barely responds: the ratio of peak
side-to-side to peak fore-aft acceleration is of order \(\sim\!40\), confirming
that even a broadband electrical event is filtered by the tower into an almost
purely side-to-side motion. This is the mechanical-band-pass picture of
\S\ref{subsec:se-selectivity}, seen for a discrete event
(Figure~\ref{fig:nt-gttrip}).

\begin{figure}[t]
  \centering
  % \includegraphics[width=0.8\linewidth]{fig/gt_trip.png}
  \caption{Discrete grid event (reduced model): a broadband platform-load step
  produces a decaying side-to-side ring-down at the tower modal frequency
  (\SI{0.234}{\hertz}), while the fore-aft mode stays essentially flat. The tower
  band-pass-filters the electrical disturbance onto its side-to-side mode.}
  \label{fig:nt-gttrip}
\end{figure}

\subsection{The frequency-support trade-off}
\label{subsec:nt-tradeoff}

The headline scenario quantifies the central tension of the thesis: the wind
turbine's frequency-support droop (\S\ref{sec:outer-droop}) helps the grid but
loads the tower. A \emph{temporary} \SI{5}{\mega\watt} load step is applied at
the main bus (smooth \SI{2}{\second} ramps, held for \SI{50}{\second}, then
removed), and the co-simulation is run twice with the identical de-loaded
operating point: once with the droop off and once with it on
(\(K_{\mathrm{droop}}=\SI{0.75}{\pu\per\hertz}\), \SI{10}{\percent} headroom).
Because the tower modes are forward-coupled (\S\ref{subsec:stm-forward}), the two
runs differ only in the droop action, so the comparison is clean.

The results are collected in Table~\ref{tab:nt-tradeoff}. On the grid side the
droop helps: it cuts the settled (plateau) frequency dip by \SI{20}{\percent}
(from \SI{0.089}{\hertz} to \SI{0.071}{\hertz}) and the transient dip by
\SI{15}{\percent}, by injecting an extra \(\approx\SI{0.008}{\pu}\) of active power
from the headroom reserve. On the mechanical side the same injection is the
penalty: the side-to-side acceleration grows by a factor \(1.72\) in both peak
and RMS. The fore-aft response stays negligible. The interaction is thus not
merely academic --- providing a genuine grid service (frequency support) directly
and measurably increases the resonance-prone tower loading, a coupling that a
turbine-level or grid-level design would need to weigh against each other.

\begin{table}[t]
  \centering
  \caption{Frequency-support trade-off (reduced model): temporary
  \SI{5}{\mega\watt} load step, droop off versus on
  (\(K_{\mathrm{droop}}=\SI{0.75}{\pu\per\hertz}\), \SI{10}{\percent} headroom).
  The droop reduces the grid-frequency dip but increases the tower side-to-side
  loading.}
  \label{tab:nt-tradeoff}
  \begin{tabular}{lccc}
    \toprule
    Quantity & Off & On & Ratio on/off \\
    \midrule
    Transient freq.\ dip [\si{\hertz}]     & 0.0897 & 0.0760 & 0.85 \\
    Plateau freq.\ dip [\si{\hertz}]       & 0.0893 & 0.0714 & 0.80 \\
    WT power, plateau [\si{\pu}]            & 0.1432 & 0.1513 & --- \\
    \midrule
    SS accel.\ peak [\si{\metre\per\second\squared}] & 0.0137 & 0.0235 & 1.72 \\
    SS accel.\ RMS  [\si{\metre\per\second\squared}] & 0.0028 & 0.0049 & 1.72 \\
    FA accel.\ peak [\si{\metre\per\second\squared}] & 0.0005 & 0.0009 & --- \\
    \bottomrule
  \end{tabular}
\end{table}

\begin{figure}[t]
  \centering
  % \includegraphics[width=0.8\linewidth]{fig/gt_trip_support.png}
  \caption{Frequency-support trade-off (reduced model): grid frequency (top), WT
  active power (middle) and tower side-to-side acceleration (bottom) for the
  temporary load step, with the droop off versus on. The droop lifts the power
  and shrinks the frequency dip, but rings the side-to-side mode
  \(\approx\!1.7\times\) harder.}
  \label{fig:nt-tradeoff}
\end{figure}

\subsection{Resonant excitation by a cyclic process load}
\label{subsec:nt-resonance}

Where the load step is broadband and only \emph{rings} the tower down, this
scenario supplies a \emph{sustained} excitation in the tower band and drives the
side-to-side mode into genuine resonance. It represents a realistic islanded-grid
situation: a large cyclic mechanical load --- a reciprocating compressor or a
beam / sucker-rod pump running near 14 strokes per minute
(\(=\SI{0.234}{\hertz}\)) --- appears as a periodic active-power oscillation at
the main bus. A \(\pm\SI{5}{\mega\watt}\) sinusoidal load at \SI{0.234}{\hertz}
is switched on at \(t=\SI{20}{\second}\), producing a grid-frequency oscillation
of \(\approx\SI{67}{\milli\hertz}\) amplitude at that frequency.

Because \SI{0.234}{\hertz} coincides with the side-to-side natural frequency, the
mode is driven resonantly through the open path \eqref{eq:se-path} and its
amplitude \emph{builds up} over roughly \(3\text{--}4\,\tau_{\mathrm{ss}}\approx
\SI{45}{\second}\) to a large sustained value, rather than decaying as in
\S\ref{subsec:nt-gttrip} (Figure~\ref{fig:nt-resonance}). The settled side-to-side
lock-in acceleration reaches
\(A_{\mathrm{ss}} = \SI{0.191}{\metre\per\second\squared}\) --- about four times
the peak of the broadband step --- while the fore-aft lock-in is only
\(A_{\mathrm{fa}} = \SI{0.0063}{\metre\per\second\squared}\), a side-to-side/
fore-aft ratio of \(\approx\!30\). The build-up versus ring-down contrast is the
defining signature of resonant, as opposed to broadband, excitation.

\begin{figure}[t]
  \centering
  % \includegraphics[width=0.8\linewidth]{fig/resonance_0234.png}
  \caption{Resonant excitation (reduced model): a \(\pm\SI{5}{\mega\watt}\)
  cyclic process load at \SI{0.234}{\hertz} (top) drives a
  \(\approx\SI{67}{\milli\hertz}\) grid-frequency oscillation (middle), and the
  side-to-side mode builds up over \(\sim\!\SI{40}{\second}\) to a sustained
  amplitude while the fore-aft mode stays flat (bottom, with the side-to-side
  amplitude envelope dashed).}
  \label{fig:nt-resonance}
\end{figure}

\subsection{The resonance curve: frequency selectivity of the coupling}
\label{subsec:nt-curve}

Sweeping the cyclic-load frequency from \SIrange{0.200}{0.300}{\hertz} and
measuring the settled side-to-side lock-in amplitude at each frequency
(\texttt{sweep\_ss\_resonance.py}) traces out the resonance curve of
Table~\ref{tab:nt-curve} and Figure~\ref{fig:nt-curve}. It is a clean, sharp
Lorentzian peak: the side-to-side amplitude climbs from
\SI{0.051}{\metre\per\second\squared} at \SI{0.200}{\hertz} to a maximum of
\SI{0.190}{\metre\per\second\squared} at exactly \SI{0.234}{\hertz} --- matching
both the calibrated modal frequency and the dedicated resonant run
(\S\ref{subsec:nt-resonance}) --- and falls away symmetrically above it. The
half-power width is only \(\approx\SI{0.026}{\hertz}\), a quality factor
\(Q\approx 9\) consistent with \(\zeta_{\mathrm{ss}}=\num{0.05}\)
(\(Q=1/2\zeta=10\)). The grid-frequency drive is essentially flat
(\(\approx\SI{67}{\milli\hertz}\)) across the whole band, so the peak is the
tower's mechanical resonance and not a feature of the excitation. The fore-aft
amplitude stays below \SI{0.0065}{\metre\per\second\squared} everywhere,
including where the load frequency crosses the (nearly coincident) fore-aft
frequency --- once more the drive path, not the frequency, is what selects the
side-to-side mode.

The practical reading is the frequency selectivity itself: a cyclic platform load
produces significant tower loading only if its fundamental or a harmonic falls
within the narrow \(\approx\SI{0.026}{\hertz}\) window around
\SI{0.234}{\hertz}; outside it the coupling is negligible. This both bounds the
risk and identifies precisely which operational frequencies to avoid.

\begin{table}[t]
  \centering
  \caption{Reduced-model resonance curve: settled side-to-side and fore-aft
  lock-in acceleration versus cyclic-load frequency (\(\pm\SI{5}{\mega\watt}\)).
  The bold row marks the side-to-side mode.}
  \label{tab:nt-curve}
  \begin{tabular}{lcc}
    \toprule
    \(f\) [\si{\hertz}] &
      \(A_{\mathrm{ss}}\) [\si{\metre\per\second\squared}] &
      \(A_{\mathrm{fa}}\) [\si{\metre\per\second\squared}] \\
    \midrule
    0.200          & 0.0507 & 0.0009 \\
    0.216          & 0.0948 & 0.0015 \\
    0.226          & 0.1538 & 0.0031 \\
    0.230          & 0.1780 & 0.0052 \\
    \textbf{0.234} & \textbf{0.1903} & \textbf{0.0064} \\
    0.238          & 0.1828 & 0.0059 \\
    0.242          & 0.1623 & 0.0038 \\
    0.248          & 0.1297 & 0.0015 \\
    0.255          & 0.1026 & 0.0011 \\
    0.280          & 0.0578 & 0.0005 \\
    0.300          & 0.0440 & 0.0003 \\
    \bottomrule
  \end{tabular}
\end{table}

\begin{figure}[t]
  \centering
  % \includegraphics[width=0.8\linewidth]{fig/ss_resonance_curve.png}
  \caption{Reduced-model resonance curve: settled side-to-side (red) and fore-aft
  (green) acceleration amplitude versus the frequency of a \(\pm\SI{5}{\mega\watt}\)
  cyclic grid load. The side-to-side response peaks sharply at \SI{0.234}{\hertz}
  (\(Q\approx 9\)); the fore-aft response stays flat, even where the load
  frequency crosses the fore-aft mode.}
  \label{fig:nt-curve}
\end{figure}

\subsection{Summary}
\label{subsec:nt-summary}

Taken together the four scenarios trace the interaction from mechanism to
operational consequence. A discrete grid event shows the tower band-pass-filtering
a broadband disturbance onto its side-to-side mode
(\S\ref{subsec:nt-gttrip}); a sustained cyclic load shows the same mode driven
into resonant build-up (\S\ref{subsec:nt-resonance}); the sweep shows the coupling
is sharply frequency-selective about \SI{0.234}{\hertz} (\S\ref{subsec:nt-curve});
and the frequency-support scenario shows that acting on the grid with the turbine
directly aggravates that tower loading (\S\ref{subsec:nt-tradeoff}). All four run
in real time on the reduced model against the full LEOGO network, complementing
the high-fidelity FMU evidence of \S\ref{sec:tower-ss-resonance} with a fast,
fully closed electrical-to-mechanical demonstration.

